% GROK REVIEW::GROK-REV-20261002-1113::2026-10-02 11:13 PT
% Independent comments. Not an LSPV ownership-region rewrite.
\documentclass[11pt]{article}
\usepackage[T1]{fontenc}
\usepackage[utf8]{inputenc}
\usepackage[margin=1in]{geometry}
\usepackage{amsmath,amssymb,booktabs}
\usepackage{hyperref}

\title{Independent review of \texttt{j15marti/pnp}\\[0.3em]
\large GROK-REV-20261002-1113}
\author{Independent reviewer}
\date{2026-10-02 11:13 PT}

\begin{document}
\maketitle

\section*{Status}
P versus NP remains open. This note records comments on the repository at tree
\texttt{afb924d5d5e8b08ad736f35bd8e7efc8655d1ed0}. It does not claim a proof of
\(P\neq NP\), and it does not promote any finite experiment to an asymptotic theorem.

\section*{Corpus reviewed}
Surveyed the repository tree, then read repository-root \texttt{main.tex},
\texttt{README.md}, \texttt{AGENTS.md}, \texttt{CODEX.md},
\texttt{comments/grok/README.md}, the 2026-10-02 10:11 review, the verified arbiter
\texttt{research/main/2026-10-02\_1056\_MAIN\_1056.tex}, and the preceding
\texttt{research/main/2026-10-02\_1012\_MAIN\_1012.tex}.
The newest root banner is the 10:58 adaptive-replay exponent guard. No matching
snapshot under \texttt{research/main/} was present for that banner.
Task files and the STRATFLAT1 case loop were not re-executed.
Ren--Williams, ECCC TR26-118 (12 July 2026), was checked against the published abstract.

\section*{What is adopted}
\begin{itemize}
\item Global status remains OPEN while any load-bearing arrow is open.
\item MAIN 1056 is the current verified arbiter. STRATFLAT1 is adopted only as a
local closure lemma, with the snapshot's own caveat that independent red-team review
was pending. Polynomial-time Skolem openings
\[
o_i=F_i(x,R,c_{\le i})
\]
flatten a polynomial-round game to \(\exists R\,\forall C\,W(x,R,C)\) with \(W\in P\),
hence to \(\Sigma_2^P\). A polynomial-size strategy with a polynomial-time evaluator
has the same shape. Under \(P=NP\), PH-PROV-CLOSE1 reacquires these cases.
\item The STRATFLAT1 control remains COMPUTATIONAL: 917{,}504 paths at
\((p,h,m,T)=(17,4,4,4)\) and 658{,}560 paths at \((29,5,5,3)\), with interactive and
flattened responses agreeing. Local ordered-basis nonuniqueness does not defeat the
canonical RREF selector.
\item SBACR1's canonical response \(F(w,s)=\operatorname{RREFBasis}(S_s(w))\) remains a
polynomial-time source evaluator. EVACR1's compositional history quotient remains a
closure for ordinary evolving restrictions.
\item The 10:58 simulation statement is adopted only after defining the missing outer
exponent. If the outer computation on length \(n\) costs \(O(n^a)\) aside from its
queries, makes at most \(n^q\) adaptive SAT queries of length at most \(n^b\), and SAT
on length \(m\) costs \(Cm^K\), direct simulation costs
\(O(n^{\max\{a,q+bK\}})\). This is an upper bound on one simulation. It is not a
no-sharing or direct-product theorem.
\item Ren--Williams TR26-118 is adopted as barrier hygiene only. The report proves a
near-maximum circuit lower bound for \(\mathrm{E}^{\mathrm{prMA}}/1\) and analyzes
\(\mathrm{P}^{\mathrm{NP}}\) with \(r\) adaptive rounds and witness length \(s\).
The abstract does not state an unrestricted SAT direct-product theorem. Transfer class:
useful structural analogy requiring a new source-matched theorem, not an exact transfer.
\item ACRTREE1, TRACECLOSE1, PH-PROV-CLOSE1, REACQ1, COVERBREAK1, EXCESS1, BITRANK1's
rejection of complex or real Gram rank, and the BOOL hostile controls remain in force.
\end{itemize}

\section*{Comments, without strengthening}
\begin{enumerate}
\item The 10:58 banner used \(a\) in \(O(n^{\max\{a,q+bK\}})\) without defining it.
The root document now states that \(a\) is the outer time exponent aside from queries.
The correction does not add a lower bound.
\item Adaptive depth does not multiply a fixed polynomial SAT exponent under direct
simulation. The cost is \(q+bK\) in the exponent, not \(K^q\). A source-matched theorem
that forbids sharing, recomputation, batching, memoization, and algebraic oracles is
still OPEN.
\item STRATFLAT1 retires polynomial alternation only in the succinct-strategy model.
A route that proved every legal winning strategy has superpolynomial representation or
evaluation complexity, and that this game is forced on every hypothetical static SAT
decider, would still be a new theorem. No such theorem is in the corpus.
\item MAIN 1056 correctly rejects QBF and DQBF proof lower bounds, and PSPACE or NEXP
strategy problems, as direct SAT transfers absent a model-preserving compiler.
\item Specialist ownership regions were not edited.
\end{enumerate}

\section*{Duplication and contradiction audit}
No contradiction was found between STRATFLAT1 and PH-PROV-CLOSE1: the flattening lemma
is the reason the fixed-PH closure applies to succinct strategies. No contradiction was
found between the 10:58 simulation bound and the continued OPEN status of a
source-matched no-sharing theorem. The 10:12 suggestion that polynomial alternation
itself might be the first surviving provenance zone is superseded by MAIN 1056, as that
snapshot already records. The 10:58 banner has no companion MAIN file; it is treated as
a root guard, not as a second arbiter.

\section*{What this review does not do}
It does not edit LSPV-1 through LSPV-4 ownership regions, does not revive retired BOOL
predictors, and does not treat prime provenance, strategy-tree cardinality, adaptive
query depth, or finite signature diversity as an unrestricted SAT-time measure.

\section*{Single next target}
On one source-legal reconstruction-compatible family, construct \(\Phi\) and the legal
slice matrices and prove or falsify \(\operatorname{rank}_{\mathbb F_2}\Phi=\omega(m)\)
by an intrinsically characteristic-two certificate. On that same family, search for
compact Skolem circuits, algebraic evaluators, decision diagrams, and source/history
quotients. A polynomial representation should be recorded as a hostile closure of the
provenance route, not as progress toward \(P\neq NP\).

% END_OF_VERIFIED_REPORT::GROK::GROK-REV-20261002-1113
\end{document}
